\documentclass[9pt,xcolor={dvipsnames}]{beamer}
\usepackage{algorithm}
\usepackage[noend]{algorithmic}
\usepackage{multirow}
\usepackage{wasysym}
\usepackage{textcomp}
\usepackage{tikz}
\usepackage{amsfonts}
\usepackage{beamerthemesplit}
\usepackage{amsmath, amssymb}
%\setbeameroption{show notes}

%\usepackage[round]{natbib}
%\renewcommand{\newblock}{}

\usepackage{mathtools}

\newcommand*{\minx}{%
    \mathmakebox[\widthof{$\max$}][r]{\min}%
}
\newcommand*{\maxx}{%
    {\max}%
}
\newcommand*{\minxx}{%
    \hphantom{\max}\mathllap{\min}%
}

\usepackage{xcolor}

\usepackage{color, colortbl}

\definecolor{LightCyan}{rgb}{0.68,1,1}
\definecolor{LightCyan2}{rgb}{0.88,1,1}
\definecolor{LightGreen}{rgb}{0.68,1,0.88}
\definecolor{LightPink}{rgb}{1,0.88,0.88}

\newcolumntype{g}{>{\columncolor{LightCyan}}c}
\newcolumntype{f}{>{\columncolor{LightCyan2}}c}
\newcolumntype{j}{>{\columncolor{LightGreen}}c}
\newcolumntype{p}{>{\columncolor{LightPink}}c}


\mode<presentation>
{
	%%% COMMENT IF YOU PREFER PREVIOUS GRAPHIC STYLE
	%%% >>>>>>>>>>>>>>>>>>>>>>>>>>>>>>>>>>>>>>>>>>>>
%	\usetheme{CambridgeUS}
%	\usecolortheme{dove}
%	\setbeamersize{text margin left=1em,text margin right=1em}
%	\defbeamertemplate{itemize item}{image}{\small\includegraphics[height=2.5ex]{world.png}}
%	\defbeamertemplate{itemize subitem}{image}{\small\includegraphics[height=2.5ex]{plane.png}}
%	\defbeamertemplate{itemize subsubitem}{image}{\small\includegraphics[height=2.5ex]{bus.png}}
%	\setbeamertemplate{itemize item}[image]
%	\setbeamertemplate{itemize subitem}[image]
%	\setbeamertemplate{itemize subsubitem}[image]
	%%% <<<<<<<<<<<<<<<<<<<<<<<<<<<<<<<<<<<<<<<<<<<<

	%%% UNCOMMENT TO GET PREVIOUS GRAPHIC STYLE
	%%% >>>>>>>>>>>>>>>>>>>>>>>>>>>>>>>>>>>>>>>
	%\usetheme{Luebeck}
	\usetheme{oro} % Beamer Theme
	\usecolortheme{oro} % Beamer Color Theme
	\usefonttheme{oro}
	%\usecolortheme{whale} % Beamer Color Theme
	\useoutertheme{oro} % Beamer Outer Theme
	%\useinnertheme{oro} % Beamer Outer Theme
	%%% <<<<<<<<<<<<<<<<<<<<<<<<<<<<<<<<<<<<<<<
}

\usepackage[T1]{fontenc}

%\titlegraphic{\includegraphics[height=3cm]{logopoli.png}}

%%% COMMENT IF YOU USE THEME "ORO"
%%% >>>>>>>>>>>>>>>>>>>>>>>>>>>>>>>>>>>>>>>
%\logo{\includegraphics[height=0.67cm]{imgs/logopolito.png}}
%\newcommand{\mostra}[1]{{\usebeamercolor[fg]{myblue}#1}}
%\newcommand{\mostraLink}[1]{{\usebeamercolor[fg]{mylink}#1}}
%%% <<<<<<<<<<<<<<<<<<<<<<<<<<<<<<<<<<<<<<<

%\usepackage{biblatex}
%\bibliography{global_references}

\title[SVCSBPP-CL-SS]{Stochastic Variable Cost and Size Bin Packing with Capacity Loss and Supplier Selection problem}
\author [Bargetto, Crainic, Perboli, Rei] %[R.~Bargetto, T.~G.~Crainic, G.~Perboli, W.~Rei]
	{Roberto~Bargetto${}^1$, Teodor~Gabriel~Crainic${}^2$, Guido~Perboli${}^{1,2}$, Walter~Rei${}^2$}
\institute[POLITO/CIRRELT]{{\small ${}^1$ Politecnico di Torino, Turin, Italy,} 	\\
{\small ${}^2$ CIRRELT, Montreal, Canada,}		\\
\vspace{1.0cm}
\includegraphics[width=5 cm]{LogoPoliCirrelt}}

\date{ISMP 2024}

\begin{document}

\usebackgroundtemplate{%
\tikz\node[opacity=0.15] {\includegraphics[height=\paperheight,width=\paperwidth]{Castello_del_valentino.jpg}};}

\setbeamertemplate{section in toc}{%
    \small\includegraphics[height=2.5ex]{world.png}  \hspace{0.3em} \inserttocsection \par}

\setbeamertemplate{subsection in toc}{
 \leavevmode\leftskip=1.5em   \small\includegraphics[height=2.5ex]{plane.png} \hspace{0.3em} \inserttocsubsection \par}

\begin{frame}[plain]
\maketitle
\end{frame}
 \usebackgroundtemplate{}

\begin{frame}
\frametitle{Outline}
\tableofcontents[pausesections]
\end{frame}

\section {Problem definition}
%%% Stochastic variable cost and size bin packing with capacity Loss and supplier selection
\begin{frame}[allowframebreaks]
	\frametitle{Problem definition}
	\framesubtitle{Supplier selection}
	\begin{itemize}
		\item The shipping business has moved from long-term contracts to a more flexible environment based on short-term contracts and ``on-spot'' transactions
		\bigskip
		\item This shift is rooted into the events that have characterized recent years (pandemic, war, etc.) and impacted logistic systems and supply chains
		\bigskip
		\item The market of shipping has been shaken and the volatility of transportation rates has become normal condition
		\begin{itemize}
			\item This asks for better decision making but leads also to better profit opportunities
		\end{itemize}
		\bigskip
		\item Shippers have now to frequently select their carrier suppliers and diversify this selection to cover the risk of undelivered capacity in reason of many possible events (mechanical failures, accidents, delays, etc.)
	\end{itemize}
\end{frame}

\begin{frame}[allowframebreaks]
	\frametitle{Problem definition}
	\framesubtitle{Supplier selection}
	\begin{itemize}
		\item Carrier companies sell storage and transportation services at
		\bigskip
		\begin{itemize}
			\item \textbf{Contract rates}
			\medskip
			\item \textbf{Spot rates}
		\end{itemize}
		\bigskip
		\item Contract rates are applied for medium term purchase agreements for securing future use of capacity and spot rates for just-in-time operations
		\bigskip
		\item \textbf{\color{purple} Spot rates} are typically much higher that \textbf{\color{blue} contract rates}
	\end{itemize}
\end{frame}

\begin{frame}[allowframebreaks]
	\frametitle{Problem definition}
	\framesubtitle{Supplier selection}
	\begin{itemize}
		\item We consider the problem of a \textbf{shipper} seeking a better profit on operations 
		\bigskip
		\item The shipper wants to identify the set of carrier companies to contract with for the next period of operations
		\bigskip
		\item The business setting of the next period is uncertain, the shipper does not known perfectly
		\bigskip
		\begin{itemize}
			\item the amount of orders will be received
			\bigskip
			\item the actual capacity the carrier companies will actually deliver if contracted
			\bigskip
			\item the spot rates applied
		\end{itemize}
	\end{itemize}
\end{frame}

\begin{frame}[allowframebreaks]
	\frametitle{Problem definition}
	\framesubtitle{Supplier selection}
	\begin{itemize}
%		\item \textbf{\color{purple} Spot rates} are typically much higher than \textbf{\color{blue} contract rates}
%		\bigskip
		\item The shipper looks for the best trade-off of capacity bought at contract rates and spot rates
		\begin{itemize}
			\footnotesize
			\item we assume also that the shipper faces a cost for capacity loss in the 2nd-stage and that the concerned suppliers pay a compensation to the shipper
		\end{itemize}
		\bigskip
		\item The shipper wants to compose a portfolio of carrier suppliers that minimizes the expected cost of operations
		\bigskip
		\item The shipper reduces the risk of uncertain factors diversifying the supplier selection
	\end{itemize}
\end{frame}

\section{Two-stage stochastic integer programming model}
\begin{frame}[allowframebreaks]
	\frametitle{Two-stage stochastic integer programming model}
	\framesubtitle{1st-stage and 2nd-stage decisions}
	\begin{itemize}
		\item The shipper's supplier selection problem is two-stage by definition
		\bigskip
		\item The shipper
		\bigskip
		\begin{itemize}
		\item selects the suppliers and buy some capacity at contract rates in the \textbf{1st stage} and
		\bigskip
		\item buy at spot rates the capacity required to ship all orders in the \textbf{2nd stage}, when complete information is available
		\end{itemize}
	\end{itemize}
\end{frame}

\begin{frame}[allowframebreaks]
	\frametitle{Two-stage stochastic integer programming model}
	\framesubtitle{Sets and notation}
	\begin{itemize}
		\small
		\item We model the supplier selection problem as an extension of the two-stage modeling of the Stochastic Variable Cost and Size Bin Packing Problem (SVCSBPP) of~\cite{crainic2024}
		\medskip
		\item In~\cite{crainic2024}, the focus is on the uncertain 2nd-stage capacity of the bins selected in the 1st-stage from a single supplier (uncertainty is due to possible mechanical failures, delays of delivery, etc.)
		\medskip
		\item A previous work, i.e.,~\cite{crainic2016logistics}, considered the problem of packing an uncertain set of items into a set of bins select in the 1st-stage to resort to the purchase of additional bins in the 2nd-stage
		\medskip
		\item Neglecting uncertainty, tactical planning problems of transportation are effectively modeled and solved as VCSBPP, e.g., last-mile delivery~\cite{perboli2021new}
		\medskip
		\item The same problem setting of~\cite{crainic2016logistics} is considered in~\cite{bruni2023machine} form the perspective of e-commerce retailers (such as Amazon or Alibaba) and forecasting machine learning techniques are used to match real-time computational requirements
		\medskip
		\item Progressive-hedging appears as the best algorithm for two-stage SVCSBPPs, e.g.,~\cite{bruni2023machine, crainic2016logistics, crainic2024}
	\end{itemize}
\end{frame}

\begin{frame}[allowframebreaks]
	\frametitle{Two-stage stochastic integer programming model}
	\framesubtitle{Sets and notation}
	\begin{itemize}
		\item Let's see the two-stage stochastic problem formulation
		\bigskip
		\item Given
		\begin{description}
			\footnotesize
			\item [$\mathcal{M}$:] the set of available suppliers
			\medskip
			\item [$T$ and $ \cal T$:] the bin types available respectively at the 1st-stage and 2nd-stage
			\medskip
			\item [$\mathcal{J}^t_m$:] the set of bins of type $ t$ that the supplier $ m$ can provide at the 1st-stage
			\medskip
			\item [$\theta_m$:] the reliability of supplier $ m$ (i.e., the percentage actually delivered on average in the past)
			\medskip
			\item [$R$ and $\nu$:] the average minimum reliability the shipper targets for selected suppliers and the minimum percentage of capacity to buy at 1st-stage from each selected supplier respectively
			\medskip
			\item [$\Omega$:] the set of 2nd-stage scenarios
			\medskip
			\item [${\cal K}(\omega)$:] the set of 2nd-stage bins available under scenario $ \omega \in \Omega$
			\medskip
			\item [$I_s(\omega)$ and $D$:] the set of order of size $ d$ under scenario $ \omega \in \Omega$ and the set of possible order sizes respectively
		\end{description}
	\end{itemize}
\end{frame}

\begin{frame}[allowframebreaks]
	\frametitle{Two-stage stochastic integer programming model}
	\framesubtitle{1st-stage variables}
	\begin{itemize}
		\item Given also
		\begin{description}
			\footnotesize
			\item [$ f_m$ and $ f^t_m$:] respectively the cost for contracting supplier $ m$ and the contract rate for bins of type $ t$ from supplier $ m$
			\medskip
			\item [$c^t_m$ and $ r^t_m$:] respectively the cost the shipper face and the compensation the supplier pays for lost capacity (we assume $c^t_m \ge r^t_m$)
			\medskip
			\item [$ V^t_m$, $ V^\tau$ and ${\cal V}_{j}(\omega)$:] respectively the volumes of 1st-stage and 2nd-stage bins according to type $ t$ or $ \tau$ and of the 1st-stage bin $ j$ under scenario $ \omega \in \Omega$
		\end{description}
		\bigskip
		\item 1st-stage variables
		\begin{description}
			\footnotesize
			\item [$y_m$:] a binary variable that takes value 1 if supplier $ m \in \mathcal{M}$ is selected
			\medskip
			\item [$y^t_{jm}$:] a binary variable that takes value 1 if bin $ j$ of type $ t \in T$ is bought from supplier $ m \in \mathcal{M}$
		\end{description}
		\bigskip
		\item and 2nd-stage variables
		\begin{description}
			\footnotesize
			\item [$x_{dj}$:] an integer variable counting the number of orders of size $ d$ to pack into bin $ j$
			\medskip
			\item [$z^\tau_{k}(\omega)$:] a binary variable that takes value 1 if 2nd-stage bin $ k$ of type $ \tau \in {\cal T}$ is used
		\end{description}
	\end{itemize}
\end{frame}

\begin{frame}[allowframebreaks]
	\frametitle{Two-stage stochastic integer programming model}
	\framesubtitle{1st-stage problem}
	We have
	\bigskip
	{
	\footnotesize
	\begin{equation}
		\label{2s:objf}
		\min_{y} \sum\limits_{m \in \cal M} {f_{m}}y_{m} +  \sum\limits_{t \in T} \sum\limits_{m \in {{\cal M}}} \sum\limits_{j \in {\cal J}_{m}^t} f_{m}^t y_{j,m}^t  + {\mathbb{E}}_{\xi}\left[Q \left(y,\xi(\omega)\right)\right] 
	\end{equation}
	\noindent subject to
	\begin{alignat}{2}
		\sum\limits_{m \in \cal {M}} y_{m}\geq{N}, & \qquad
		\label{eq3}\\
		\sum\limits_{t \in T} \sum\limits_{j \in {\cal J}^t_{m}} V_{m}^{t} \ y_{j,m}^{t} \geq \nu \sum\limits_{t \in T} \sum\limits_{m \in {\cal M}} \sum\limits_{j \in {\cal J}^t_{m}} V_{m}^{t} \ y_{j,m}^{t}, & \qquad \forall m \in {\cal M}
		\label{eq3.2}\\
		\sum\limits_{m \in {{\cal M}}} \ \theta_{m} \ y_{m}  \geq {R} \sum\limits_{m \in {\cal M}} \ y_{m}, & \qquad \label{eq4}\\
		\sum\limits_{t \in T} \sum\limits_{j \in {\cal J}^t_{m}} y_{j,m}^{t} \leq y_{m}, & \qquad \forall m \in \cal M, \label{eq5}\\
		y_{j,m}^t \in \{0, 1\}, & \qquad \forall t \in T, m \in {\cal M}, j \in {\cal J}^t_{m}, \label{eq6}\\
		y_{m} \in \{0, 1\}, & \qquad \forall m \in \cal M \label{eq7},
	\end{alignat}
	}
\end{frame}

\begin{frame}[allowframebreaks]
	\frametitle{Two-stage stochastic integer programming model}
	\framesubtitle{2nd-stage problem}
	where
	{
	\footnotesize
	\begin{equation}
		\begin{split}
			Q \left(y,\xi (\omega)\right) = & \min_{z(\omega), x(\omega)} 
			\sum\limits_{\tau \in \cal T}\sum\limits_{k \in \cal K^{\tau}(\omega)} g^{\tau}(\omega) z_k^{\tau}(\omega) \; + \\
			& + \sum\limits_{t \in T} \sum\limits_{m \in {\cal M}} \sum\limits_{j \in {\cal J}^t_{m}} (c_{m}^t-r_{m}^t)(V^{t} - {\cal V}_{j}(\omega)) y_{j,m}^t \label{eq8}
		\end{split}
	\end{equation}
	\noindent {subject to}
	\begin{alignat}{2}
		\sum\limits_{j \in \cal J \cup {\cal K} (\omega)} x_{d,j}(\omega) = |I_d(\omega)|, & \qquad \forall d \in D, \label{eq9}\\
		\sum\limits_{d \in D} s \cdot x_{d,j}(\omega) \leq {\cal V}_{j}(\omega) y_{j,m}^t, & \qquad \forall \: t \in T, m\in{{\cal M}}, j \in {\cal J}^t_{m}, \label{eq10}\\
		\sum\limits_{d \in D} s \cdot x_{d,k}(\omega) \leq V^\tau z^{\tau}_k(\omega), & \qquad \forall \: \tau \in {\cal T}, k \in {\cal K}^{\tau}(\omega), \label{eq11}\\
		x_{d,j}(\omega) \in \mathbb{Z}^+, & \qquad \forall \: d \in D, j \in {\cal J}, \label{eq12}\\
		x_{d,k}(\omega) \in \mathbb{Z}^+, & \qquad \forall \: d \in D, k \in \cal K(\omega), \label{eq13}\\
		z_k^{\tau}(\omega) \in \{0, 1\}, & \qquad \forall \: \tau \in {\cal T}, k \in {\cal K}^{\tau}(\omega). \label{eq14}
	\end{alignat}
	}
\end{frame}


\section{Row aggregation lower bound}
\begin{frame}[allowframebreaks]
	\frametitle{Row aggregation of the recourse problem}
	\begin{itemize}
		\item Given the total volume of orders $\Upsilon (\omega) = \sum\limits_{i \in I(\omega)} v_i$ under scenario $ \omega \in \Omega$
		\bigskip
		\item We can aggregate constraints~\eqref{eq10} and~\eqref{eq11} over bins, bin types and suppliers, and consequently discard constraints~\eqref{eq9}
		\bigskip
		\item And we obtain recourse problem
		{
			\footnotesize
			\begin{equation*}
				Q \left(y,\xi (\omega)\right)
			\end{equation*}
			\noindent {subject to}
			\begin{alignat}{2}
				\Upsilon (\omega) \leq \sum\limits_{\substack{ m \in {\cal M}, \\ t \in T, \\ j \in {\cal J}^t_{m}}} {\cal V}_{j}(\omega) y_{j,m}^t + \sum\limits_{\substack{\tau \in {\cal T}, \\ k \in {\cal K}^{\tau,s}}} V^\tau z^{\tau}_k(\omega), & \qquad \forall \: \omega \in \Omega, \label{eq10lb}\\
				z_k^{\tau}(\omega) \in \{0, 1\}, & \qquad \forall \: \tau \in {\cal T}, k \in {\cal K}^{\tau}(\omega), \label{eq14lb}
			\end{alignat}
		}
	\end{itemize}
\end{frame}

\begin{frame}[allowframebreaks]
	\frametitle{Row aggregation of the recourse problem}
	\begin{itemize}
		\item By solving problem
		{
		\footnotesize
		\begin{equation*}
		\min_{y} \sum\limits_{m \in \cal M} {f_{m}}y_{m} +  \sum\limits_{t \in T} \sum\limits_{m \in {{\cal M}}} \sum\limits_{j \in {\cal J}_{m}^t} f_{m}^t y_{j,m}^t  + {\mathbb{E}}_{\xi}\left[Q \left(y,\xi(\omega)\right)\right] 
		\end{equation*}
		\noindent {subject to}
		\begin{equation*}
			\eqref{eq3},~\eqref{eq3.2},~\eqref{eq4},~\eqref{eq5},~\eqref{eq6},~\text{and}~\eqref{eq10lb}
		\end{equation*}
		\begin{alignat*}{2}
		y_{j,m}^t \in \{0, 1\}, & \qquad \forall t \in T, m \in {\cal M}, j \in {\cal J}^t_{m}, \\
		y_{m} \in \{0, 1\}, & \qquad \forall m \in \cal M, \\
		z_k^{\tau}(\omega) \in \{0, 1\}, & \qquad \forall \: \tau \in {\cal T}, k \in {\cal K}^{\tau}(\omega),
		\end{alignat*}
		}
		we obtain a \textbf{valid lower bound} of the problem
		\bigskip
		\item Note that we discarded all scenario-dependent packing variables $ x_{dj}(\omega)$
		\bigskip
		\item The obtained solution of 1st-stage variables $ y_m$ and $y^t_{jm}$ can then be evaluated in the original problem for a set of scenarios $ \tilde{\Omega} \subset \Omega$
		\begin{itemize}
			\item This step provides an \textbf{upper bound} for the \textbf{recourse problem} (RP) evaluation
		\end{itemize}
	\end{itemize}
\end{frame}

\section {Numerical results}
\begin{frame}[allowframebreaks]
	\frametitle{Evaluation deterministic-equivalent ILPs}
	\begin{itemize}
		\item We evaluate the deterministic-equivalent ILPs of the two-stage problem
		\begin{itemize}
			\item~\eqref{2s:objf}-\eqref{eq14} and
			\item that of the row aggregation lower bound
		\end{itemize}
		\bigskip
		\item The deterministic-equivalent ILP of~\eqref{2s:objf}-\eqref{eq14} can be solved only for a limited number of scenarios (5, 10, 20)
		\begin{itemize}
			\footnotesize
			\item Let $ S$ be the set of generated scenario approximating $\Omega$
		\end{itemize}
		\bigskip
		\item We generate sets of instances of small size 100-200 orders per scenario and 3 suppliers
		\bigskip
		\item To speed up the model resolution we implemented a procedure for the complete enumeration of supplier selections (this is possible due to small number of suppliers of an instance)
	\end{itemize}
\end{frame}

\begin{frame}[allowframebreaks]
	\frametitle{Instance generation}
	Generated instances reflect the following (relevant) characteristics
	\begin{itemize}
		\footnotesize
		\medskip
		\item Three suppliers, $|{\cal M}| = 3$
		\medskip
		\item Select at least two suppliers, $N = 2$, and a minimum volume of bin capacity of 20\%, $\nu = 0.2$
		\medskip
		\item With minimum ave. reliability of 80\%, $R = 0.8$
		\medskip
		\item Reliability of each supplier $\theta_{m}$ drawn from $U(0.7, 0.9)$
		\medskip
		\item Nr. of items per scenario drawn from $U(100, 200)$
		\medskip
		\item Three 1st-stage bin types with volumes $ V_t \in \{ \, 50, \, 100, \, 150 \, \}$
		\medskip
		\item Nr. of 2st-stage bins per bin type $ | {\cal J}^t_m | = \lceil \frac{\text{MTVIS}}{V^t} \rceil $; MTVIS is the maximum total volume of items over the generated scenarios $ S$
		\medskip
		\item Nr. of 1st-stage bins per type $ | {\cal K}^t_m |$ drawn from $ U_{\text{int}}( 0.1 | {\cal J}^t_m |, 0.5 | {\cal J}^t_m |)$
	\end{itemize}
\end{frame}

\begin{frame}[allowframebreaks]
	\frametitle{Instance generation}
	\begin{itemize}
%		\footnotesize
		\item Reliability of supplier $ m$ under scenario $ s \in S$ is computed as $ \tilde{\theta}^s_m = \theta_m + \varepsilon$ with $  \varepsilon $ drawn from
		\begin{itemize}
			\footnotesize
			\item  $ U(-0.1, 0.1)$, \textbf{low variance} (LV), or
			\item  $ U(-0.3, 0.3)$, \textbf{high variance} (HV)
		\end{itemize}
		\bigskip
		\item Aggregated capacity loss of supplier $ m$ under scenario $ s \in S$ is computed as ${\cal L}^s_m = (1 - \tilde{\theta}^s_m) \sum_{i \in I_s} \nu^s_i$ 
		\bigskip
		\item 2nd-stage loss ${\cal V}^s_j$ of 1st-stage bin $ j$ under scenario $ s \in S$, is computed according to two cases
		\begin{description}
			\footnotesize
			\item [\textit{Uniform}:] ${\cal V}^s_j = V^t_m - l^s_m$ with $ l^s_m = \frac{{\cal L}^s_m}{|{\cal J}^t_m|} $
			\item [\textit{Localized}:] ${\cal V}^s_j$ is $0$ or $ V^t_m $ such that $ \sum_{j \in {\cal J}^t} {\cal V}^s_j + {\cal L}^s_m = V^t_m |{\cal J}^t| $
		\end{description}
	\end{itemize}		
\end{frame}

\begin{frame}[allowframebreaks]
	\frametitle{Instance generation}
	\begin{itemize}
		\item 1st-stage costs
		\begin{itemize}
			\item Contract cost $ f_m = \zeta_m \cdot \nu \cdot \max_{s \in S} \sum_{i \in I_s} \nu^s_i $ with $\zeta_m \sim U(0.1, 0.5)$
			\item Contract rates $f^t_m = V^t (1 + \gamma^t_m)$ with $\gamma^t_m \sim U(-0.25,0.25)$
		\end{itemize}
		\bigskip
		\item 2nd-stage costs	
		\begin{itemize}
			\item Capacity loss cost $c^t_m = 1.25 \frac{f^t_m}{V^t}$ and compensation $r^t_m = \frac{f_m^t}{V^t}$
			\item Spot rates $g^t_m = f^t_m (1 + \beta^t_m)$ with 
			\begin{itemize}
				\bigskip
				\item $\beta^t_m \sim U(0.25, 0.5)$, \textbf{small spread} (SS), or
				\item $\beta^t_m \sim U(0.5, 0.75)$, \textbf{large spread} (LS)
			\end{itemize}
		\end{itemize}
	\end{itemize}		
\end{frame}

\begin{frame}[allowframebreaks]
	\frametitle{Numerical experiments}
	\framesubtitle{Expected Value of Perfect Information, EVPI, and Value of Stochastic Solution, VSS}
	\begin{itemize}
		\footnotesize
		\item 20 instances for each set, 20 or 120 scenarios for each instance
		\bigskip
		\item All four combination of low/high variance of supplier reliability (LV/HV) and small/large spread of spot rates (SS/LS) w.r.t. contract rates
		\bigskip
		\item For the WS and EEV computation, the model w/ packing vars does not always solve to optimality the single-scenario problem, but the gap is limited (< 1\%)
		\begin{itemize}
			\item if not optimal, the best bound is taken
			\item the same for the RP objective function evaluation
		\end{itemize}
		\bigskip
		\item All the experiments are run as single thread processes on a PC equipped with a 12th Gen Intel Core i9-12900 3.8GHz, with 32GB of RAM and running Ubuntu 22.04 LTS
		\bigskip
		\item Time-limits are
		\begin{itemize}
			\item 600 sec. for each supplier selection of the complete enumeration procedure, when there is no scenario separation (i.e., EV and RP computation)
			\item  30 sec. for each supplier selection in the case of scenario separation (i.e., WS and EEV computation)
		\end{itemize}
	\end{itemize}
\end{frame}

\begin{frame}[allowframebreaks]
	\frametitle{Numerical experiments}
	\framesubtitle{Expected Value of Perfect Information, EVPI, and Value of Stochastic Solution, VSS}
	Model with packing variables (PV) and model row aggregation (RA)
%	\begin{itemize}
%		\item Comparison of model with packing vars vs. row aggregation model
%		\bigskip
%		\item 20 instances for each set, 20 scenarios for each instance
%		\bigskip
%		\item Instances with uniform loss of bin capacity
%	\end{itemize}
%	\bigskip
	\begin{table}[h!]
		\centering
		\footnotesize
		\begin{tabular}{c c c j f c}
			\hline
			Model & 
			Loss type & 
			Instance set & 
			EVPI & 
			VSS & 
			RP gap \\
			\hline
			PV & U & LV-SS & 2.14\% & \textbf{8.13\%} & 1.75\%  \\
			PV & U & LV-LS & 5.18\% & 3.58\% & 3.98\% \\
			PV & U & HV-SS & 4.24\% & \textbf{8.58\%} & 1.43\% \\
			PV & U & HV-LS & \textbf{7.21\%} & 5.36\% & 4.87\% \\
			\hline
			PV & L & LR-SS & 4.98\% & 3.58\% & 4.25\% \\
			PV & L & LR-LS & \textbf{7.64\%} & 1.98\% & 3.81\% \\
			PV & L & HR-SS & 4.81\% & 3.83\% & 3.14\% \\
			PV & L & HR-LS & \textbf{8.43\%} & 2.22\% & 3.08\% \\
			\hline
			RA & U & LV-SS & 2.35\% & \textbf{7.84\%} & 0.18\% \\
			RA & U & LV-LS & 5.60\% & 3.15\% & 0.24\% \\
			RA & U & HV-SS & 4.42\% & \textbf{8.31\%} & 0.11\% \\
			RA & U & HV-LS & \textbf{7.68\%} & 4.79\% & 0.33\% \\
			\hline
			RA & L & LR-SS & 5.60\% & 2.70\% & 0.00\%  \\
			RA & L & LR-LS & \textbf{8.33\%} & 1.21\% & 0.00\%  \\
			RA & L & HR-SS & 5.52\% & 2.97\% & 0.00\% \\
			RA & L & HR-LS & \textbf{9.09\%} & 1.49\% & 0.00\% \\
			\hline
		\end{tabular}
		\caption{EVPI and VSS evaluation, 20 instances for each set, 20 scenarios; uniform loss (U) or localized loss (L)}
		\label{tab:optres}
	\end{table}
%	\begin{itemize}
%		\item Instances with uniform loss seems slightly harder to solve
%		\item For instances with localized loss, the EVPI\% always dominates VSS\%
%	\end{itemize}
\end{frame}

\begin{frame}[allowframebreaks]
	\frametitle{Numerical experiments}
	\framesubtitle{Expected Value of Perfect Information, EVPI, and Value of Stochastic Solution, VSS}
	\begin{itemize}
		\footnotesize
		\item We evaluate the quality of the RP solution we obtain if we inject the RA model solution of 1st-stage variable into model PV
		\begin{itemize}
			\item Note that by doing so we have separation of scenarios and a much faster evaluation
		\end{itemize}
		\medskip 
		\item The evaluation is done by comparing this RP solution and that of model PV as the gap w.r.t. the best RP bound (PV or RA)
	\end{itemize}
	\begin{table}[h!]
		\centering
		\footnotesize
		\begin{tabular}{c c j c f c f}
			\hline
%			 & PV & RA &  & PV &  & RA &  \\
			Loss &
			Set & 
			Best RP bound & 
			RP of PV &
			Gap & 
			RP injected sol. & 
			Gap w/ best \\
			\hline
			U & LV-SS & 3094.7 & 3146.6 & \textbf{1.68\%} & 3157.2 & 2.02\% \\
			U & LV-LS & 3370.2 & 3503.9 & 3.97\% & 3402.5 & 0.96\% \\
			U & HV-SS & 3165.2 & 3210.4 & \textbf{1.43\%} & 3232.5 & 2.13\% \\
			U & HV-LS & 3225.9 & 3384.7 & 4.92\% & 3260.7 & \textbf{1.08\%} \\
			\hline
			L & LV-SS & 2782.9 & 2899.7 & 4.20\% & 2810.2 & \textbf{0.98\%} \\
			L & LV-LS & 2948.1 & 3060.8 & 3.82\% & 2976.9 & \textbf{0.98\%} \\
			L & HV-SS & 2815.0 & 2903.4 & 3.14\% & 2841.7 & \textbf{0.95\%} \\
			L & HV-LS & 2827.4 & 2915.1 & 3.10\% & 2851.9 & \textbf{0.87\%} \\
			\hline
		\end{tabular}
		\caption{RP evaluation, model with packing variables (PV) \textit{vs.} injecting row aggregation (RA) 1-stage solution}
		\label{tab:optres4}
	\end{table}

\end{frame}

\begin{frame}[allowframebreaks]
	\frametitle{Numerical experiments}
	\framesubtitle{Expected Value of Perfect Information, EVPI, and Value of Stochastic Solution, VSS}	
	\begin{itemize}
%		\footnotesize
		\item We evaluate also the row aggregation model for a larger set of scenarios; 120 scenarios
	\end{itemize}
	\bigskip
	\begin{table}[h!]
		\centering
		\footnotesize
		\begin{tabular}{r r j f r}
			\hline
			Loss type & Instance set & EVPI & VSS & RP gap \\
			\hline
			U & LVRLSR & 2.17\% & \textbf{9.47\%} & 0.15\% \\
			U & LVRHSR & 5.73\% & 5.52\% & 0.40\% \\
			U & HVRLSR & 4.74\% & \textbf{8.92\%} & 0.18\% \\
			U & HVRHSR & \textbf{7.57\%} & 6.45\% & 0.47\% \\
			\hline
			L & HVRHSR & \textbf{8.84\%} & 1.80\% & 0.00\% \\
			L & HVRLSR & \textbf{5.87\%} & 2.52\% & 0.00\% \\
			L & LVRHSR & \textbf{8.59\%} & 1.76\% & 0.00\% \\
			L & LVRLSR & \textbf{5.80\%} & 2.98\% & 0.00\% \\
			\hline
		\end{tabular}
		\caption{EVPI and VSS evaluation, 120 scenarios; uniform loss (U) localized loss (L)}
		\label{tab:optres3}
	\end{table}
%	\begin{itemize}
%		\item Instances with loss uniform over bins seems slightly harder to solve
%		\item For instances with localized loss, the EVPI\% always dominates VSS\%
%	\end{itemize}
\end{frame}

\begin{frame}[allowframebreaks]
	\frametitle{Numerical experiments}
	\framesubtitle{Some considerations}	
	\begin{itemize}
		\item Instances with uniform loss of bin capacity seems to be slightly harder to solve
		\bigskip
		\item For tested instances, characterized by a small amount of items, the EVPI and VSS indicators suggest a relevant impact of stochastic optimization for the problem
		\bigskip
		\item The large number of ILP variables due to the several suppliers to select makes the problem hard to solve, meta-heuristic approaches reducing the solution space can be effective
		\bigskip
		\item The tested procedure of 1st-stage variables solution injection suggests as promising stochastic solution algorithms allowing separation of scenarios (e.g., progressive-hedging)
	\end{itemize}
\end{frame}

\section {What's next?}
\begin{frame}[allowframebreaks]
	\frametitle{Looking the way forward}
	\begin{itemize}
		\item Devise a meta-heuristic solution approach for the selection of suppliers 
		\bigskip
		\item Devise a pseudo-exact progressive-hedging algorithm for solving the recourse problem
	\end{itemize}
\end{frame}


\begin{frame}[shrink]
\frametitle{References}
\begin{thebibliography}{9}
\bibitem {bruni2023machine} Maria Elena Bruni, Edoardo Fadda, Stanislav Fedorov, Guido Perboli
\textbf{ A machine learning optimization approach for last-mile delivery and third-party logistics},
Computers \& Operations Research, Vol. 157, 2023

\bibitem {crainic2016logistics} Teodor Gabriel Crainic, Luca Gobbato, Guido Perboli,  Walter Rei
\textbf{ Logistics capacity planning: A stochastic bin packing formulation and a progressive hedging meta-heuristic},
	European Journal of Operational Research, Vol. 253, No. 2, 2016, pp. 404-417

\bibitem {crainic2024} Teodor Gabriel Crainic, Guido Perboli,  Walter Rei, Rosano Mariangela, Veronica Lerma
\textbf{ Capacity planning with uncertainty on contract fulfillment},
	European Journal of Operational Research, Vol. 314, No. 1, 2024, pp. 152-175

\bibitem {perboli2021new} Guido Perboli, Luce Brotcorne, Maria Elena Bruni, Mariangela Rosano
\textbf{ Transportation Research Part E: Logistics and Transportation Review},
	European Journal of Operational Research, Vol. 145, 2021


\end{thebibliography}
\end{frame}

%\usebackgroundtemplate{%
%\tikz\node[opacity=1.0] {\includegraphics[height=\paperheight]{thinking.png}};}
%\begin{frame}[plain]
%\end{frame}
%
%\usebackgroundtemplate{}

\end{document}